\documentclass[
	12pt,
	a4paper,
	bibtotoc,
	cleardoubleempty, 
	idxtotoc,
	ngerman,
	openright
	final,
	listof=nochaptergap,
	]{scrbook}
\usepackage{cmap}
\usepackage[T1]{fontenc}
\usepackage[utf8]{inputenc}

% ##################################################
% Unterstuetzung fuer die deutsche Sprache
% ##################################################
%\usepackage{ngerman}
\usepackage[ngerman]{babel}
\usepackage{graphicx}
\usepackage{wrapfig}

% ##################################################
% Dokumentvariablen
% ##################################################

% Persoenliche Daten
\newcommand{\docNachname}{Mustermann}
\newcommand{\docVorname}{Max}
\newcommand{\docStrasse}{Straße}
\newcommand{\docOrt}{Ort}
\newcommand{\docPlz}{123456}
\newcommand{\docEmail}{student@hs-furtwangen.de}
\newcommand{\docMatrikelnummer}{000000}

% Dokumentdaten
\newcommand{\docTitle}{TITEL}
%\newcommand{\docUntertitle}{} % Kein Untertitel
\newcommand{\docUntertitle}{UNTERTITEL}
% Arten der Arbeit: Bachelorthesis, Masterthesis, Seminararbeit, Diplomarbeit
\newcommand{\docArtDerArbeit}{ART DER ARBEIT}
%Studiengaenge: Allgemeine Informatik Bachelor, Computer Networking Bachelor,
% Software-Produktmanagement Bachelor, Advanced Computer Scinece Master
\newcommand{\docStudiengang}{STUDIENGANG}
\newcommand{\docAbgabedatum}{11.11.2011}
\newcommand{\docErsterReferent}{ERSTER REFERENT}
%\newcommand{\docZweiterReferent}{-} % Wenn es nur einen Betreuer gibt
\newcommand{\docZweiterReferent}{ZWEITER REFERENT}

% ##################################################
% Allgemeine Pakete
% ##################################################

% Abbildungen einbinden
\usepackage{graphicx}

% Zusaetsliche Sonderzeichen
\usepackage{dingbat}

% Symbole Haken und X [OPTIONAL]
%\usepackage{pifont}
%\newcommand{\cmark}{\ding{51}}
%\newcommand{\xmark}{\ding{55}}

% Farben
\usepackage{color}
\usepackage[usenames,dvipsnames,svgnames,table]{xcolor}

% Maskierung von URLs und Dateipfaden
\usepackage[hyphens]{url}

% Deutsche Anfuehrungszeichen
\usepackage[babel, german=quotes]{csquotes}

% Pakte zur Index-Erstellung (Schlagwortverzeichnis)
\usepackage{index}
\makeindex

% Ipsum Lorem
% Paket wird nur für das Beispiel gebraucht und kann gelöscht werden
\usepackage{lipsum}

% ##################################################
% Seitenformatierung
% ##################################################
\usepackage[
	portrait,
	bindingoffset=1.5cm,
	inner=2.5cm,
	outer=2.5cm,
	top=3cm,
	bottom=2cm,
	%showframe, %Aktivieren um Seitengrenzen anzuzeigen
	%includeheadfoot
	]{geometry}

% ##################################################
% Kopf- und Fusszeile
% ##################################################

\usepackage{fancyhdr}

\pagestyle{fancy}
\fancyhf{}
\fancyhead[EL,OR]{\sffamily\thepage}
\fancyhead[ER,OL]{\sffamily\nouppercase{\leftmark}}

\fancypagestyle{headings}{}

\fancypagestyle{plain}{}

\fancypagestyle{empty}{
  \fancyhf{}
  \renewcommand{\headrulewidth}{0pt}
}

%Speichert \chaptermark in \oldchaptermark damit 
% es für die Anhänge zurückgesetzt werden kann
\let\oldchaptermark\chaptermark

%Kein "Kapitel # NAME" in der Kopfzeile
\renewcommand{\chaptermark}[1]{
	\markboth{#1}{}
   	\markboth{\thechapter.\ #1}{}
}

% ##################################################
% Schriften
% ##################################################

% Stdandardschrift festlegen
\renewcommand{\familydefault}{\sfdefault}

% Standard Zeilenabstand: 1,5 zeilig
\usepackage{setspace}
\onehalfspacing 

% Schriftgroessen festlegen
\addtokomafont{chapter}{\sffamily\large\bfseries} 
\addtokomafont{section}{\sffamily\normalsize\bfseries} 
\addtokomafont{subsection}{\sffamily\normalsize\mdseries} 
\addtokomafont{caption}{\sffamily\normalsize\mdseries} 

%Einrücken von Absätzen deaktivieren
\setlength{\parindent}{0pt}

%Zeilenabstand bei abstätzen
\usepackage{parskip}

% ##################################################
% Inhaltsverzeichnis / Allgemeine Verzeichniseinstellungen
% ##################################################

\usepackage{tocloft}

% Punkte auch bei Kapiteln
\renewcommand{\cftchapdotsep}{3}
\renewcommand{\cftdotsep}{3}

% Schriftart und -groesse im Inhaltsverzeichnis anpassen
\renewcommand{\cftchapfont}{\sffamily\normalsize}
\renewcommand{\cftsecfont}{\sffamily\normalsize}
\renewcommand{\cftsubsecfont}{\sffamily\normalsize}
\renewcommand{\cftchappagefont}{\sffamily\normalsize}
\renewcommand{\cftsecpagefont}{\sffamily\normalsize}
\renewcommand{\cftsubsecpagefont}{\sffamily\normalsize}

%Zeilenabstand in den Verzeichnissen einstellen
\setlength{\cftparskip}{.5\baselineskip}
\setlength{\cftbeforechapskip}{.1\baselineskip}

%Einrücken von Absätzen deaktivieren
%\setlength{\parindent}{0pt}

%Zeilenabstand bei abstätzen
\usepackage{parskip}

% ##################################################
% Abbildungsverzeichnis und Abbildungen
% ##################################################

\usepackage{caption}

\usepackage{wrapfig}

% Nummerierung von Abbildungen
\renewcommand{\thefigure}{\arabic{figure}}
\usepackage{chngcntr}
\counterwithout{figure}{chapter}

% Abbildungsverzeichnis anpassen
\renewcommand{\cftfigpresnum}{Abbildung }
\renewcommand{\cftfigaftersnum}{:}

% Breite des Nummerierungsbereiches [Abbildung 1:]
\newlength{\figureLength}
\settowidth{\figureLength}{\bfseries\cftfigpresnum\cftfigaftersnum}
\addtolength{\figureLength}{2mm} %extra offset
\setlength{\cftfignumwidth}{\figureLength}
\setlength{\cftfigindent}{0cm}

% Schriftart anpassen
\renewcommand\cftfigfont{\sffamily}
\renewcommand\cftfigpagefont{\sffamily}

%standardpfad anpassen
\graphicspath{ {../src/content/pictures/} }

% ##################################################
% Tabellenverzeichnis und Tabellen
% ##################################################

% Nummerierung von Tabellen
\renewcommand{\thetable}{\arabic{table}}
\counterwithout{table}{chapter}

% Tabellenverzeichnis anpassen
\renewcommand{\cfttabpresnum}{Tabelle }
\renewcommand{\cfttabaftersnum}{:}

% Breite des Nummerierungsbereiches [Abbildung 1:]
\newlength{\tableLength}
\settowidth{\tableLength}{\bfseries\cfttabpresnum\cfttabaftersnum}
\addtolength{\tableLength}{3mm} %extra offset
\setlength{\cfttabnumwidth}{\tableLength}
\setlength{\cfttabindent}{0cm}

%Schriftart anpassen
\renewcommand\cfttabfont{\sffamily}
\renewcommand\cfttabpagefont{\sffamily}

% Unterdrueckung von vertikalen Linien
\usepackage{booktabs}

%Multi row für spezifische zellen
\usepackage{multirow}

%Additional table package
\usepackage{tabu}

% ##################################################
% Listings (Quellcode)
% ##################################################

\usepackage{listings}

%use typewriter font which supports bold characters
\usepackage{beramono}

\definecolor{codegreen}{rgb}{0,0.6,0}
\definecolor{codegray}{rgb}{0.5,0.5,0.5}
\definecolor{codepurple}{rgb}{0.5,0,0.33}
\definecolor{codepurblue}{rgb}{0.16,0.0,1.0}
\definecolor{backcolour}{rgb}{0.95,0.95,0.92}

\lstdefinestyle{codestyle}{
    backgroundcolor=\color{backcolour},   
    commentstyle=\color{codegreen},
    keywordstyle=\bfseries\color{codepurple},
    numberstyle=\tiny\color{codegray},
    stringstyle=\color{codepurblue},
    basicstyle=\scriptsize\ttfamily,
    breakatwhitespace=false,         
    breaklines=true,                 
    captionpos=b,                    
    keepspaces=true,                 
    numbers=left,                     
    numbersep=5pt,                 
    showspaces=false,                
    showstringspaces=false,
    showtabs=false,                  
    tabsize=2
}

\lstset{style=codestyle}

%Code auschnitt importieren aus datei
%\mylisting{from}{to}{language}{file}{descr}{path}
\newcommand{\mylisting}[6]{
\lstinputlisting[language=#3,
				firstnumber=#1,
				firstline=#1,
				lastline=#2,
				caption={#4, #5}, 
				label={implementation_listing_#4_#5}]
				{#6}
}

% ##################################################
% Appendix
% ##################################################

%Calc packet für berechnungen
\usepackage{calc}

%Appendix paket, setzen der flags für das TOC
\usepackage[toc,title,titletoc]{appendix} 

%Umbenennen der überschrift für die Anhänge 
\renewcommand{\appendixtocname}{Anhänge}

%Befehl für einen neuen Bericht und die erste seite als bild
\newcommand{\appendixsection}[2]{
\section{#1}
\appendixsingle{#2}
}

%Befehl für einzelne seite als bild eingefasst, damit überschrift und kopfzeile
% bestehen bleibt. 
\newcommand{\appendixsingle}[1]{
\vspace{-10cm}
\vfill
\mbox{}\hspace{-1.5cm}\includegraphics[width=\linewidth+3cm]{#1}\hspace{-1.5cm}\mbox{}
\vspace{-10cm}
\vfill
\mbox{}
}

%Datenträger Tabelle
\definecolor{lightgray}{gray}{0.85}
\definecolor{ultralightgray}{gray}{0.95}
\definecolor{mygray}{gray}{0.70}

% ##################################################
% Theoreme
% ##################################################
  	
% Umgebung fuer Beispiele
\newtheorem{beispiel}{Beispiel}

% Umgebung fuer These
\newtheorem{these}{These}

% Umgebung fuer Definitionen
\newtheorem{definition}{Definition}
  	
% ##################################################
% Literaturverzeichnis
% ##################################################

\usepackage{bibgerm}

% ##################################################
% Abkuerzungsverzeichnis
% ##################################################

\usepackage[printonlyused]{acronym}

% ##################################################
% PDF / Dokumenteninternelinks
% ##################################################

\usepackage[
	colorlinks=false,
   	linkcolor=black,
   	citecolor=black,
  	filecolor=black,
	urlcolor=black,
    bookmarks=true,
    bookmarksopen=true,
    bookmarksopenlevel=3,
    bookmarksnumbered,
    plainpages=false,
    pdfpagelabels=true,
    hyperfootnotes,
    pdftitle ={\docTitle},
    pdfauthor={\docVorname~\docNachname},
    pdfcreator={\docVorname~\docNachname}]{hyperref}

\begin{document}

\setcounter{secnumdepth}{3}

% Titelblatt
\begin{titlepage}
\pagestyle{empty}

% ##################################################
% HFU-Logo einbinden
% ##################################################
\begin{flushright}
\begin{figure}[ht]
\flushright
\includegraphics[height=3cm]{content/pictures/hfu.jpg}
\end{figure}
\end{flushright}

% ##################################################
% Titel
% ##################################################
\begin{center}
{\fontsize{18}{22} \selectfont \docArtDerArbeit}\\[5mm]
{\fontsize{18}{22} \selectfont im Studiengang} \\[5mm]
{\fontsize{18}{22} \selectfont \docStudiengang}\\
\vspace{1cm}
\begin{onehalfspace}
{\fontsize{22}{26} \selectfont \textbf{\docTitle}}\\[5mm]
{\fontsize{18}{22} \selectfont \docUntertitle}


\end{onehalfspace}
\end{center}

% ##################################################
% Zusatzinformationen
% ##################################################
\vfill
\begin{center}
\begin{tabular}{lcl}
Referent  		&:& \docErsterReferent 	\\ \\
Koreferent 		&:& \docZweiterReferent \\ \\	
Vorgelegt am 	&:& \docAbgabedatum 	\\ \\
Vorgelegt von 	&:& \docVorname~\docNachname\\
				& & Matrikelnummer: \docMatrikelnummer\\
				& & \docStrasse,~\docPlz~\docOrt	\\
				& & \docEmail			
\end{tabular}
\end{center}
\end{titlepage}
\cleardoubleemptypage

\frontmatter
\pagenumbering{Roman}

% Abstract
\chapter*{Abstract\markboth{Abstract}{}}
\addcontentsline{toc}{chapter}{Abstract}

[Englisches Abstract (100-120 Worte)]

[Deutsches Abstract (100-120 Worte)]
\cleardoubleemptypage

% Inhaltsverzeichnis
\phantomsection
\addcontentsline{toc}{chapter}{Inhaltsverzeichnis}
\tableofcontents
\cleardoubleemptypage

% Abbildungsverzeichnis einbinden und ins Inhaltsverzeichnis
% WORKAROUND: tocloft und KOMA funktionieren zusammen nicht
% korrekt\phantomsection
\phantomsection 
\addcontentsline{toc}{chapter}{\listfigurename} 
\listoffigures
\cleardoubleemptypage

% Tabellenverzeichnis einbinden und ins Inhaltsverzeichnis
% WORKAROUND: tocloft und KOMA funktionieren zusammen nicht
% korrekt\phantomsection
\phantomsection
\addcontentsline{toc}{chapter}{\listtablename}
\listoftables
\cleardoubleemptypage

% Quellcodeverzeichnis einbinden und ins Inhaltsverzeichnis
\phantomsection
\addcontentsline{toc}{chapter}{Quellcodeverzeichnis}

%Define listing
\makeatletter
\begingroup\let\newcounter\@gobble\let\setcounter\@gobbletwo
  \globaldefs\@ne \let\c@loldepth\@ne
  \newlistof{listings}{lol}{\lstlistlistingname}
\endgroup
\let\l@lstlisting\l@listings
\makeatother
\setlength{\cftlistingsindent}{0em}
\renewcommand{\cftlistingsafterpnum}{\vskipspt}
\renewcommand*{\cftlistingspresnum}{\lstlistingname~}
\settowidth{\cftlistingsnumwidth}{\cftlistingspresnum}
\renewcommand{\lstlistlistingname}{Quellcodeverzeichnis}
% Tabellenverzeichnis anpassen
\renewcommand{\lstlistingname}{Codeauschnitt}
\renewcommand{\cftlistingsaftersnum}{:}
% Breite des Nummerierungsbereiches [Codeauschnitt 1:]
\newlength{\codeLength}
\settowidth{\codeLength}{\bfseries\lstlistingname\cftlistingsaftersnum}
\addtolength{\codeLength}{5mm}
\setlength{\cftlistingsnumwidth}{\codeLength}
\lstlistoflistings
\cleardoubleemptypage

% Abkürzungsverzeichnis
\chapter*{Abkürzungsverzeichnis\markboth{Abkürzungsverzeichnis}{}}
\addcontentsline{toc}{chapter}{Abkürzungsverzeichnis}

\begin{acronym}
\acro{HFU}{Hochschule Furtwangen University}
\end{acronym}

\mainmatter

\chapter{Einleitung}

\section{Strukturen}

\subsection{asdasd}
asdfasdf

\subsubsection{Subsubsection}
asdf

\paragraph{Paragraph}
asdf

\section{Beispielabbildungen}

\lipsum[10]

\begin{wrapfigure}{rt}{8cm}
\caption{Bildüberschrift}
\centering
\includegraphics[width=0.3\textwidth]{content/pictures/hfu}
Quelle: \cite{s11wasml}
\label{pic:bild2}
\end{wrapfigure}

\lipsum[10]

\begin{figure}
\caption{Bildüberschrift}
\includegraphics[width=1\textwidth]{content/pictures/hfu}
Quelle: \cite{s11wasml}
\label{pic:bild1}
\end{figure}

\section{Beispieltabelle}

\begin{table}
\caption{Tabellenüberschrift}
\center
\footnotesize
\begin{tabular}{lll}
\toprule
Head1 & Head2 & Head3 \\
\midrule
Val1 & Val2 & Val3 \\
Val4 & Val5 & Val6 \\
\bottomrule
\end{tabular}
\end{table}

\section{Listings}

\begin{lstlisting}[language=java, caption=Hallo Welt in Java]
public class HalloWelt 
{
	public static void main(String[] args) 
	{
		System.out.println("Hallo Welt!");
	}
}
\end{lstlisting}

%\mylisting{from}{to}{language}{file}{descr}{path}
\mylisting{1}{9}{java}{Java}{Alles}{content/code/example.java}
\mylisting{6}{7}{java}{Java}{Ausschnitt}{content/code/example.java}

\section{Abkürzungen}

Im Abkürzungsverzeichnis stehende Abkürzungen können in Langform (\ac{HFU}) oder in Kurzform (\acs{HFU}) angegeben werden.

\section{Abgesetztes wörtliches Zitat}

\lipsum[10]

\begin{quote}
\textit{\enquote{Eingerücktes wörtliches Zitat.}}\cite[S. 14ff]{s11wasml}
\end{quote}

\lipsum[10]

\section{Theoreme}

\lipsum[2]

\begin{beispiel}
Beispieltext\dots
\end{beispiel}
 
\lipsum[2]

\begin{these}
These\dots
\end{these}
 
\lipsum[2]
 
\begin{definition}
Unter dem Begriff \dots verstehen wir \dots
\end{definition}

\lipsum[2]
\chapter{Grundlagen}
\begin{figure}
\caption{Bildüberschrift}
\includegraphics[width=1\textwidth]{content/pictures/hfu}
\label{pic:bild3}
\end{figure}
\chapter{[Eigene Kapitel]}
\chapter{Ausblick} 
\chapter{Fazit}

% Schalgwortverzeichnis (Index)
%\printindex

% Literaturverzeichnis
\singlespacing
\bibliographystyle{alphadin}
\bibliography{bibtex}

% Eidesstattliche Erklärung
\chapter*{Eidesstattliche Erklärung\markboth{Eidesstattliche Erklärung}{}}
% Eintrag in das Inhaltsverzeichnis 
\addcontentsline{toc}{chapter}{Eidesstattliche Erklärung}

Ich versichere, dass ich die vorstehende Arbeit selbständig verfasst und hierzu
keine anderen als die angegebenen Hilfsmittel verwendet habe. Alle Stellen der Arbeit die 
wörtlich oder sinngemäß aus fremden Quellen entnommen wurden, sind als solche kenntlich gemacht.
\\
\\
Die Arbeit wurde bisher in gleicher oder ähnlicher Form in keinem anderen
Studiengang als Prüfungsleistung vorgelegt oder an anderer Stelle
veröffentlicht.
\\
\\
Ich bin mir bewusst, dass eine falsche Erklärung rechtliche Folgen haben kann.

\vspace*{1.5cm} \par
\line(1,0){200} \par
\docOrt, \docAbgabedatum ~~\docVorname~\docNachname

%Zurücksetzen \chaptermark
\let\chaptermark\oldchaptermark

% Hier können Anhaenge angefuegt werden
\begin{appendices}
\chapter{Monatsberichte}
Nun folgen die Monatsberichte
\newpage
\appendixsection{Monatsbericht Februar}{../report/report.pdf}
\appendixsection{Monatsbericht Januar}{../report/report.pdf}
\end{appendices}
\end{document}      